% Editable English-sentence patch record for Bethe 1928 rev22.
% IMPORTANT: This file records ONLY prose changes. Mathematical notation is
% intentionally not re-typeset here; the compiled 91-page document inherits all
% inline/display mathematics, Miller indices, tables, and figures from rev20.

% PDF p.4 / original journal pp.55-56
\newcommand{\RevP55SelectiveReflection}{%
For a given incident de Broglie wavelength, one generally observes only continuous scattering in all directions; only at certain discrete incident wavelengths are the electrons preferentially diffracted into one particular direction, or into several symmetry-related directions.%
}
\newcommand{\RevP55FootnoteFourTail}{%
Following Davisson and Germer, we denote the latter by the Miller indices of the simplest crystal face whose normal lies in the plane of reflection determined by the incident and reflected beams.%
}
\newcommand{\RevP56LatticeSentence}{%
From this phenomenon of selective reflection one may conclude that the electron waves are diffracted by the crystal lattice of nickel.%
}

% PDF p.7 / original journal p.58
\newcommand{\RevP58ThreeDimensionalAction}{%
the other hand, which depends on the three-dimensional action of the atoms in the interior of the crystal, is of course influenced by a refractive index attributed to the interior of the crystal.%
}

% PDF p.13 / original journal pp.62-63
\newcommand{\RevP63SecondOrderBeam}{%
Thus the longest second-order beam reflected by the (440) plane in the (111) azimuth should emerge grazing, but is not observed.%
}
\newcommand{\RevP63ReliableObservations}{%
This is probably due to the fact that in the second order the intensity falls off much more rapidly with decreasing angle than in the first, and already at the point where reliable observations first become possible, the scattered waves of the uppermost atomic layer have exactly the opposite phase to those of the second.%
}

% PDF p.14 / original journal p.64
\newcommand{\RevP64DispersionSentence}{%
The vanishing of their determinant yields, as usual, the dispersion equation, from which two different values result for the as yet unknown wave number.%
}

% PDF p.16 / original journal pp.66-67
\newcommand{\RevP66FieldPoints}{%
This equation cannot be satisfied simultaneously at all field points.%
}
\newcommand{\RevP67Instead}{%
Instead, it turns out%
}

% PDF p.18 / original journal pp.68-69
\newcommand{\RevP69RemainingFreedom}{%
atoms. The only remaining freedom is the choice of the propagation vector and the amplitude of the primary wave, but only so long as the crystal is really infinitely extended and nowhere comes into contact with another medium.%
}
\newcommand{\RevP69SurfaceConstants}{%
The surface phenomena treated in the following section determine the constants of the ``primary beam''.%
}

% PDF p.36 / original journal pp.86-87
\newcommand{\RevParticleFluxSentence}{%
In this wave description, the electron number density is proportional to the squared modulus of the de Broglie-wave amplitude. We denote the particle flux of the primary beam in vacuum by the quantity defined in Eq.~(56).%
}

% PDF p.48 / original journal pp.96-97
\newcommand{\RevWeakBeamReaction}{%
lattice point from the sphere of propagation that the intensity of the ``weak beams'' becomes vanishingly small, no ``reaction'' need be taken into account. The same occurs for electrons of much smaller wavelengths.%
}

% PDF p.49 / original journal pp.97-98
\newcommand{\RevDirectCoupling}{%
can no longer couple directly to the primary wave. For in the dynamical equation%
}

% PDF p.50 / original journal pp.99-100
\newcommand{\RevIndirectReflection}{%
For high electron velocities the ``indirect reflection'' due to weak-beam back-coupling is of course no longer relevant. At low velocities it may be observed, provided a conducting crystal with a suitable structure is found.%
}

% PDF p.76 / original journal pp.128-129, Summary
\newcommand{\RevSummaryStitch}{%
circumstances. Section 9 treats the case, realized by Davisson and Germer, of reflection in several directions symmetric with respect to the incident beam. The mean crystal potential acting on the electrons, responsible for the refractive index, is found on purely theoretical grounds, on the basis of wave-mechanical assumptions about the structure of the crystal atoms, to be nearly of the same magnitude as follows from the experiments.%
}
